\documentclass[noauthor]{ximera}
 
\title{Week 12}
\author{Math 1193 Design Team}
 
\begin{document}
\begin{abstract}
    Solutions for Week 12.
\end{abstract}
\maketitle
 
\begin{problem}
Which of the following expressions are equal to \(abc\)? There are multiple
answers.
\begin{enumerate}
    \item \(\left((abc)^{-1}\right)^{-1}\)
    \item \((abc)^{-1}a^{-1}b^{-1}c^{-1}\)
    \item \(\frac{bc}{\frac{1}{a^{-1}}}\)
    \item \((abc + 1)^2 - 1\)
    \item \(\frac{\sqrt{(ab)^2}}{c^{-2}}\)
    \item \(\sqrt{a}\sqrt{b}\sqrt{c}a^{1/2}b^{1/2}c^{1/2}\)
    \item \(\frac{(a^{2}b^{2})^{1/2}}{c^{-1}}\)
    \item \(\frac{\sqrt{(-a)^2}((cb)^2 - 1)}{cb + 1} + a\)
    \item \(\left(\frac{a^2b^2}{c^{-2}}\right)^{1/2}\)
\end{enumerate}
\begin{explanation}
\begin{enumerate}
    \item \(\left((abc)^{-1}\right)^{-1} = (abc)^{(-1)(-1)} = (abc)^1 = abc\), so yes.
    \item \((abc)^{-1}a^{-1}b^{-1}c^{-1} = (abc)^{-1}(abc)^{-1} = (abc)^{-2}\), so no.
    \item \(\frac{bc}{\frac{1}{a^{-1}}} = \frac{bc}{a}\), so no.
    \item \((abc + 1)^2 - 1 = (abc)^2 + 2abc + 1 - 1 = (abc)^2 + 2abc\), so no.
    \item \(\frac{\sqrt{(ab)^2}}{c^{-2}} = ((ab)^2)^{1/2}c^2 = abc^2\), so no.
    \item \(\sqrt{a}\sqrt{b}\sqrt{c}a^{1/2}b^{1/2}c^{1/2} = a^{1/2}b^{1/2}c^{1/2}a^{1/2}b^{1/2}c^{1/2} = a^{1/2 + 1/2}b^{1/2 + 1/2}c^{1/2 + 1/2} = a^1b^1c^1 = abc\), so yes.
    \item \(\frac{(a^{2}b^{2})^{1/2}}{c^{-1}} = \frac{(a^2)^{1/2}(b^2)^{1/2}}{c^{-1}} = a^{2 \cdot 1/2}b^{2 \cdot 1/2}c^1 = abc\), so yes.
    \item \(\frac{\sqrt{(-a)^2}((cb)^2 - 1)}{cb + 1} + a = \frac{(a^2)^{1/2}(cb - 1)(cb + 1)}{cb + 1} + a = a(cb - 1) + a = acb - a + a = abc\), so yes.
    \item \(\left(\frac{a^2b^2}{c^{-2}}\right)^{1/2} = \frac{(a^2)^{1/2}(b^2)^{1/2}}{(c^{-2})^{1/2}} = \frac{ab}{c^{-1}} = abc\), so yes.
\end{enumerate}
\end{explanation}
\end{problem}
 
\begin{problem}
Let \(f\) be defined by \(f(x) = \frac{3x}{x + 1}\).
 
\begin{enumerate}
\item \(f\) is invertible. How can you tell?
\item Is \(g(x) = \frac{3x}{x^2 + 1}\) invertible?
\item Find the domain of \(f\).
\item Evaluate the following expressions:
\begin{enumerate}
    \item \(f^{-1}(f(2))\)
    \item \(f^{-1}(f(-1))\)
\end{enumerate}
\item Find a formula for \(f^{-1}(x)\).
\item Find the range of \(f\).
\end{enumerate}
 
\begin{explanation}
\begin{enumerate}
    \item The graph of \(f\) passes the horizontal line test, meaning each output corresponds to only one input, meaning that \(f\) is invertible.
    \item The graph of \(g\) does not pass the horizontal line test, meaning there are some outputs that correspond to more than one input, meaning \(g\) is not invertible.
    \item Since \(x = -1\) results in the denominator being 0, \(f(-1)\) is undefined. Since \(f(x)\) is defined for every other \(x\)-value,
    the domain of \(f\) is \((-\infty, -1) \cup (-1, \infty)\), all real numbers except for \(-1\).
    \item \begin{enumerate}
        \item By the definition of inverses, since \(f(2)\) is defined, this is 2.
        \item Since \(f(-1)\) is undefined, so is \(f^{-1}(f(-1))\).
    \end{enumerate}
    \item To find the inverse, we switch \(x\) and \(y\) and solve for \(y\).
    \begin{align*}
        x & = \frac{3y}{y + 1} \\
        x(y + 1) & = 3y \\
        xy + x & = 3y \\
        x  = 3y - xy \\
        x & = y(3 - x) \\
        \frac{x}{3 - x} & = y \\
        f^{-1}(x) = \frac{x}{3 - x}
    \end{align*}
    \item The range of \(f\) is the domain of \(f^{-1}\). Since \(f^{-1}\) is only undefined at \(x = 3\), its domain is \((-\infty, 3) \cup (3 \infty)\).
    This means that the range of \(f\) is also  \((-\infty, 3) \cup (3 \infty)\).
\end{enumerate}
\end{explanation}
\end{problem}
 
 
 
\begin{problem}
\includegraphics[width=10cm]{base3.png}
\begin{enumerate}
    \item Above is the graph of \(g(x) = 3^x\). Label at least
    three points on the graph.
    \item Reflect your graph across the line \(y = x\). To do this, swap the \(x\) and \(y\) coordinates of each
    point on the graph. Call your new function \(f\).
    \begin{enumerate}
        \item What is the relationship between \(g\) and \(f\)?
        \item What are three points on the graph of \(f\)?
        \item What is another name for \(f\)? There are two possible answers:
        see if you can come up with both!
    \end{enumerate}
    \item Is the asympote to the graph of \(f\) vertical or horizontal?
    What is its equation?
    \item What are the domain and range of \(f\)? If you get
    stuck, review the relationship between the domains and
    ranges of inverse functions.
     
\end{enumerate}
 
\begin{explanation}
\begin{enumerate}
    \item Some examples of points are \((0, 1)\), \((1, 3)\), \((2, 9)\), or \((-1, 1/3)\).
    Notice that when you plug in big negative numbers like \(x = -12\), \(3^x\) is close to 0, but not equal to 0.
    This means that \((-12, 0)\) is not on the graph, but \((-12, 1/3^{12})\) is.
    \item \begin{enumerate}
        \item \(g\) and \(f\) are inverses, since switching \(x\) and \(y\) is what taking the inverse function does.
        \item Some examples of points are \((1, 0)\), \((3, 1)\), \((9, 2)\), or \((1/3, -1)\) or \((1/3^{12}, -12)\).
        \item Two other names for \(f\) are \(g^{-1}\) and \(\log_3\).
    \end{enumerate}
    \item Since \(g\) has a horizontal asymptote at \(y = 0\), \(g^{-1}\) has a vertical asymptote at \(x = 0\). This can be 
    explained by remembering that we switched \(x\) and \(y\).
    \item The domain of \(g\) is \((-\infty, \infty)\), and the range of \(g\) is \((0, \infty)\), based on the graph.
    Since the domain of the inverse is the range of the original function, the domain of \(f\) is \((0, \infty)\).
    Since the range of the inverse is the domain of the original function, the range of \(f\) is \((-\infty, \infty)\).
\end{enumerate}
\end{explanation}
\end{problem}
 
\begin{problem}
    \begin{enumerate}
        \item Help! I made a table containing some points on the graph of \(f(x) = x^2\), but my inconsiderate roommate erased parts of the table. Please fill it in for me. My professor is very picky about exact answers, so don't round! Are there any entries that could have more than one option?
        \begin{center}
\begin{tabular}{|c|c|}
 \(x\) & \(y\) \\ \hline
 1 & 1 \\ 
 \(-2\) & 4 \\
 2 & 4 \\
 \(\sqrt{2}\) & 2 \\
 0 & 0\\
 \(-9\) & 81 \\
 \(-\sqrt{3}\) & 3 \\
 \(\sqrt{2}\) & 2\\
\end{tabular}
\end{center}
        \item Same thing, but for \(g(x) = e^x\).
                \begin{center}
\begin{tabular}{|c|c|}
 \(x\) & \(y\) \\ \hline
 1 &  \(e\)\\ 
 2 & \(e^2\) \\
 1 & \(e\) \\
 0 & 1 \\
 \(-9\) & \(\frac{1}{e^9}\)\\
 \(\frac{1}{2}\)& \(\sqrt{e}\) \\
  \(-2\)& \(\frac{1}{e^2}\) \\
  \(-\frac{1}{3}\) & \(\frac{1}{\sqrt[3]{e}}\)\\
\end{tabular}
\end{center}
    \end{enumerate}
\begin{explanation}
    Notice that some cells in the top have multiple answers. To see why, consider the graph \(y = x^2\).
    Every \(y\)-value except 0 corresponds to 2 \(x\)-values. Since \(x^2\) is not invertible, you have to make choices when writing down which \(x\) value.
\end{explanation}
\end{problem}
 
\begin{problem}
Aaliyah and Tom are having a friendly disagreement. Aaliyah thinks
\(\log(15) = \log(3) + \log(5)\) and Tom thinks
\(\log(15) = \log(3) \cdot \log(5)\).
 
Aaliyah's work is as follows:
    \begin{align*}
        \log(15) & = \log(3 \cdot 5) \\
        & = \log(10^{\log(3)} \cdot 10^{\log(5)}) \\
        & = \log(10^{\log(3) + \log(5)}) \\
        & = \log(3) + \log(5)
    \end{align*}
 
Tom's work is as follows:
    \begin{align*}
        \log(15) & = \log(3 \cdot 5) \\
        & = \log(10^{\log(3)} \cdot 10^{\log(5)}) \\
        & = \log(10^{\log(3) \cdot \log(5)}) \\
        & = \log(3) \cdot \log(5)
    \end{align*}
 
 
Whose work is correct? Why? How could you check?

\begin{explanation}
    Aaliyah's work is correct. 
    
    Recall that \(b^{x}\cdot b^y = b^{x + y}\). To see this, remember that
    \(b^x\) is the product of \(x\) copies of \(b\), and \(b^y\) is the product of \(y\) copies of \(b\).
    When multiplied, you are taking the product of \(x + y\) copies of \(b\).

    Therefore, \(10^{\log(3)} \cdot 10^{\log(5)} = 10^{\log(3) + \log(5)}\). Even though, \(\log(3)\) and \(\log(5)\) are unfamiliar
    numbers, the properties of exponents still hold. 
\end{explanation}
\end{problem}
 
\begin{problem}
When I was learning logarithms, there was no way to put \(\log_2(5)\)
into my calculator, since there was only a \(\log\) button and a
\(\ln\) button, without the option to directly input a base other
than 10 or \(e\). In this exercise, we'll learn how my generation
dealt with this problem.
 
Make sure to write down your work at each step.
 
\begin{enumerate}
    \item Write the equation \(\log_2(5) = x\) in exponential form.
    \item Now take the \textbf{natural log} of both sides of your
    exponential equation.
    \item Next, we'll use a property of logarithms we'll learn soon:
    \(\ln(2^x) = x \ln(2)\). Use this to rewrite one side of your equation.
    \item Finally, solve for \(x\).
    \item How did I put \(\log_2(5)\) into my calculator? What about
    \(\log_7(19)\)?
\end{enumerate}

\begin{explanation}
    Here's the work:
    \begin{align*}
        \log_2(5) & = x \\
        2^x & = 5 \\
        \ln(2^x) & = \ln(5) \\
        x\ln(2) & = \ln(5) \\
        x & = \frac{\ln(5)}{\ln(2)}
    \end{align*}
    Therefore, instead of plugging \(\log_2(5)\) into my calculator, I plugged in \(\ln(5)/\ln(2)\). Same for \(\ln(19)/\ln(7)\).
\end{explanation}
\end{problem}
 
\begin{problem}
\begin{enumerate}
    \item Sky was asked to simplify \(\ln(e^{-2})\) and said the answer was
    \(-2\). Are they correct? Why or why not? Write down your reasoning.
    \item Sky was asked to simplify \(e^{\ln(-2)}\) and said the answer was
    \(-2\). Are they correct? Why or why not? Write down your reasoning.
\end{enumerate}
\begin{explanation}
    \begin{enumerate}
        \item Since \(e^{-2}\) is defined and \(\ln\) and \(e^x\) are inverses, the expression simplifies to \(-2\).
        \item Since \(\ln(-2)\) is undefined, the whole expression is also undefined.
    \end{enumerate}
\end{explanation}
\end{problem}
 
\end{document}